% GENERATED FILE. Edit canonical lessons, then run npm run generate:books.
% canonical-source-hash: fnv1a64:be5d851dbe78fdb9
% canonical-lessons: AR-C442-mukhayyam, AR-C442-muntaja, AR-C442-shalih, AR-C442-ashala, AR-C442-sanawbar

\chapter{Camps and resorts}
\label{ch:ar-camps-and-resorts}
\hlchaptermodality{442}

\begin{chapteropening}
\textbf{By the end of this chapter:} \emph{I can say a campsite, a resort, a chalet, to set alight and pine trees.}

Complete the last lesson of chapter 442: I can say a campsite, a resort, a chalet, to set alight and pine trees.
\end{chapteropening}

\section[\texorpdfstring{\ar{مخيم} (mukhayyam)}{مخيم (mukhayyam)}]{\ar{مخيم} (mukhayyam) — a campsite; a camp}

\label{lesson:AR-C442-mukhayyam}



\begin{quote}
Before the new one: say the Arabic for to depart from, then the Arabic for to vacate.
\end{quote}

\subsection*{You'll want to know: \ar{مخيم}}
\textbf{\ar{مخيم}} — \emph{mukhayyam} — \textquotedblleft{}a campsite; a camp\textquotedblright{}.

\begin{quote}
\ar{كان} \ar{المخيم} \ar{بجانب} \ar{البحيرة}.
\end{quote}

\begin{quote}
\emph{kāna al-mukhayyam bi-jānib al-buḥayra.}
\end{quote}

\begin{quote}
The campsite was beside the lake.
\end{quote}

\begin{cousinweb}
\textbf{\ar{مخيم}} (\emph{mukhayyam}) is built on the native Arabic root \emph{kh-y-m}, \textquotedblleft{}a tent\textquotedblright{} (Lane). It is the Form II passive participle: a place set with tents.
\end{cousinweb}

\begin{grammarlens}[title={What you've built}]
One more word for getting around.
\end{grammarlens}

\subsection*{Your turn}
\begin{itemize}
\raggedright
  \item \emph{Say it:} \emph{mukhayyam}
  \item \emph{Say it:} \emph{mukhayyam}, once more
  \item \emph{Say it:} \emph{kāna al-mukhayyam bi-jānib al-buḥayra.} — The campsite was beside the lake.
  \item \emph{Recall:} say the Arabic for to depart from, then the Arabic for to vacate, then say \emph{mukhayyam} again
\end{itemize}

\begin{tcolorbox}[breakable,colback=teal!4,colframe=teal!35!black,title={Before you move on}]
What does \ar{مخيم} mean? (\textquotedblleft{}A campsite; a camp\textquotedblright{}.) Say it once more.
\end{tcolorbox}
\section[\texorpdfstring{\ar{منتجع} (muntajaʿ)}{منتجع (muntajaʿ)}]{\ar{منتجع} (muntajaʿ) — a resort}

\label{lesson:AR-C442-muntaja}



\begin{quote}
Before the new one: say the Arabic for to vacate, then the Arabic for a campsite.
\end{quote}

\subsection*{You'll want to know: \ar{منتجع}}
\textbf{\ar{منتجع}} — \emph{muntajaʿ} — \textquotedblleft{}a resort\textquotedblright{}.

\begin{quote}
\ar{قضينا} \ar{العطلة} \ar{في} \ar{منتجع} \ar{على} \ar{البحر}.
\end{quote}

\begin{quote}
\emph{qaḍaynā al-ʿuṭla fī muntajaʿ ʿalā al-baḥr.}
\end{quote}

\begin{quote}
We spent the holiday at a resort by the sea.
\end{quote}

\begin{cousinweb}
\textbf{\ar{منتجع}} (\emph{muntajaʿ}) is built on the native Arabic root \emph{n-j-ʿ}, \textquotedblleft{}to go in search of pasture\textquotedblright{} (Lane). It is the place-noun of Form VIII: once a place people went to for fresh pasture.
\end{cousinweb}

\begin{grammarlens}[title={What you've built}]
One more word for getting around.
\end{grammarlens}

\subsection*{Your turn}
\begin{itemize}
\raggedright
  \item \emph{Say it:} \emph{muntajaʿ}
  \item \emph{Say it:} \emph{muntajaʿ}, once more
  \item \emph{Say it:} \emph{qaḍaynā al-ʿuṭla fī muntajaʿ ʿalā al-baḥr.} — We spent the holiday at a resort by the sea.
  \item \emph{Recall:} say the Arabic for to vacate, then the Arabic for a campsite, then say \emph{muntajaʿ} again
\end{itemize}

\begin{tcolorbox}[breakable,colback=teal!4,colframe=teal!35!black,title={Before you move on}]
What does \ar{منتجع} mean? (\textquotedblleft{}A resort\textquotedblright{}.) Say it once more.
\end{tcolorbox}
\section[\texorpdfstring{\ar{شاليه} (shālīh)}{شاليه (shālīh)}]{\ar{شاليه} (shālīh) — a chalet}

\label{lesson:AR-C442-shalih}



\begin{quote}
Before the new one: say the Arabic for a campsite, then the Arabic for a resort.
\end{quote}

\subsection*{You'll want to know: \ar{شاليه}}
\textbf{\ar{شاليه}} — \emph{shālīh} — \textquotedblleft{}a chalet\textquotedblright{}.

\begin{quote}
\ar{استأجرنا} \ar{شاليه} \ar{صغيرا} \ar{على} \ar{البحر}.
\end{quote}

\begin{quote}
\emph{istaʾjarnā shālīh ṣaghīran ʿalā al-baḥr.}
\end{quote}

\begin{quote}
We rented a small chalet by the sea.
\end{quote}

\begin{cousinweb}
\textbf{\ar{شاليه}} (\emph{shālīh}) is usually traced to French \emph{chalet}.
\end{cousinweb}

\begin{grammarlens}[title={What you've built}]
One more word for getting around.
\end{grammarlens}

\subsection*{Your turn}
\begin{itemize}
\raggedright
  \item \emph{Say it:} \emph{shālīh}
  \item \emph{Say it:} \emph{shālīh}, once more
  \item \emph{Say it:} \emph{istaʾjarnā shālīh ṣaghīran ʿalā al-baḥr.} — We rented a small chalet by the sea.
  \item \emph{Recall:} say the Arabic for a campsite, then the Arabic for a resort, then say \emph{shālīh} again
\end{itemize}

\begin{tcolorbox}[breakable,colback=teal!4,colframe=teal!35!black,title={Before you move on}]
What does \ar{شاليه} mean? (\textquotedblleft{}A chalet\textquotedblright{}.) Say it once more.
\end{tcolorbox}
\section[\texorpdfstring{\ar{أشعل} (ashʿala)}{أشعل (ashʿala)}]{\ar{أشعل} (ashʿala) — to set alight; to light}

\label{lesson:AR-C442-ashala}



\begin{quote}
Before the new one: say the Arabic for a resort, then the Arabic for a chalet.
\end{quote}

\subsection*{You'll want to know: \ar{أشعل}}
\textbf{\ar{أشعل}} — \emph{ashʿala} — \textquotedblleft{}to set alight; to light\textquotedblright{}.

\begin{quote}
\ar{أشعل} \ar{أبي} \ar{النار} \ar{أمام} \ar{الخيمة}.
\end{quote}

\begin{quote}
\emph{ashʿala abī an-nār amāma al-khayma.}
\end{quote}

\begin{quote}
My father lit the fire in front of the tent.
\end{quote}

Present: \ar{يشعل}, \emph{yushʿil}.

It also means to switch on a lamp.

\begin{cousinweb}
\textbf{\ar{أشعل}} (\emph{ashʿala}) is built on the native Arabic root \emph{sh-ʿ-l}, \textquotedblleft{}to blaze, to flame\textquotedblright{} (Lane). It is Form IV, \textquotedblleft{}to make blaze\textquotedblright{}.
\end{cousinweb}

\begin{grammarlens}[title={What you've built}]
One more word for getting around.
\end{grammarlens}

\subsection*{Your turn}
\begin{itemize}
\raggedright
  \item \emph{Say it:} \emph{ashʿala}
  \item \emph{Say it:} \emph{ashʿala}, once more
  \item \emph{Say it:} \emph{ashʿala abī an-nār amāma al-khayma.} — My father lit the fire in front of the tent.
  \item \emph{Recall:} say the Arabic for a resort, then the Arabic for a chalet, then say \emph{ashʿala} again
\end{itemize}

\begin{tcolorbox}[breakable,colback=teal!4,colframe=teal!35!black,title={Before you move on}]
What does \ar{أشعل} mean? (\textquotedblleft{}To set alight; to light\textquotedblright{}.) Say it once more.
\end{tcolorbox}
\section[\texorpdfstring{\ar{صنوبر} (ṣanawbar)}{صنوبر (ṣanawbar)}]{\ar{صنوبر} (ṣanawbar) — pine trees; a pine}

\label{lesson:AR-C442-sanawbar}



\begin{quote}
Before the new one: say the Arabic for a chalet, then the Arabic for to set alight.
\end{quote}

\subsection*{You'll want to know: \ar{صنوبر}}
\textbf{\ar{صنوبر}} — \emph{ṣanawbar} — \textquotedblleft{}pine trees; a pine\textquotedblright{}.

\begin{quote}
\ar{نصبنا} \ar{الخيمة} \ar{تحت} \ar{أشجار} \ar{الصنوبر}.
\end{quote}

\begin{quote}
\emph{naṣabnā al-khayma taḥta ashjār aṣ-ṣanawbar.}
\end{quote}

\begin{quote}
We pitched the tent under the pine trees.
\end{quote}

\begin{cousinweb}
\textbf{\ar{صنوبر}} (\emph{ṣanawbar}) names the pine; the origin is not settled.
\end{cousinweb}

\begin{grammarlens}[title={What you've built}]
One more word for getting around.
\end{grammarlens}

\subsection*{Your turn}
\begin{itemize}
\raggedright
  \item \emph{Say it:} \emph{ṣanawbar}
  \item \emph{Say it:} \emph{ṣanawbar}, once more
  \item \emph{Say it:} \emph{naṣabnā al-khayma taḥta ashjār aṣ-ṣanawbar.} — We pitched the tent under the pine trees.
  \item \emph{Recall:} say the Arabic for a chalet, then the Arabic for to set alight, then say \emph{ṣanawbar} again
\end{itemize}

\begin{tcolorbox}[breakable,colback=teal!4,colframe=teal!35!black,title={Before you move on}]
What does \ar{صنوبر} mean? (\textquotedblleft{}Pine trees; a pine\textquotedblright{}.) Say it once more.
\end{tcolorbox}
